\input{../tex-inputs/latex-leseansicht-vorspann}

\section[Arthur Schnitzler, Hugo von Hofmannsthal, Gisela und Markus Hajek an Gustav Schwarzkopf, 17. 8. 1898]{L04507 Arthur Schnitzler, Hugo von Hofmannsthal, Gisela und Markus Hajek an
               Gustav Schwarzkopf, 17. 8. 1898}
\nopagebreak\mylabel{L04507v}
\rehead{ }\normalsize\beginnumbering\briefempfaengerindex{Schwarzkopf, Gustav@\textsc{Schwarzkopf, Gustav}!zzzHajek, Markus@\emph{von Markus Hajek}!1898-08-171@{17. 8. 1898}|(be}\briefempfaengerindex{Schwarzkopf, Gustav@\textsc{Schwarzkopf, Gustav}!zzzHajek, Gisela@\emph{von Gisela Hajek}!1898-08-171@{17. 8. 1898}|(be}\briefempfaengerindex{Schwarzkopf, Gustav@\textsc{Schwarzkopf, Gustav}!zzzHofmannsthal, Hugo von@\emph{von Hugo von Hofmannsthal}!1898-08-171@{17. 8. 1898}|(be}\briefempfaengerindex{Schwarzkopf, Gustav@\textsc{Schwarzkopf, Gustav}!zzzSchnitzler, Arthur@\emph{von Arthur Schnitzler}!1898-08-171@{17. 8. 1898}|(be}
\toendnotes[C]{\smallbreak\pagebreak[2]}
\correspDesc{Versand  durch Arthur Schnitzler, Hugo von Hofmannsthal, Gisela Hajek, Markus Hajek am 17. 8. 1898 in Ouchy
\newline{}Übermittlung  am 17. 8. 1898 in Ouchy
\newline{}Erhalt  durch Gustav Schwarzkopf im Zeitraum [18. 8. 1898
                  – 22. 8. 1898?] in Wien}\toendnotes[C]{\smallbreak}
\Standort{DLA, A:Schnitzler, HS.1985.1.1897.}
\physDesc{Bildpostkarte, 346 Zeichen
\newline{}Handschrift Arthur Schnitzler: Bleistift, deutsche Kurrent
\newline{}Handschrift Hugo von Hofmannsthal: Bleistift, deutsche Kurrent
\newline{}Handschrift Gisela Hajek: Bleistift, deutsche Kurrent
\newline{}Handschrift Markus Hajek: Bleistift, lateinische Kurrent
\newline{}Versand: Stempel: »\nobreak{}\oindex{Ouchy@\textbf{Ouchy}|pwk}Ouchy, 17. VIII. 98, \textcolor{gray}{2}\nobreak{}«.  }\toendnotes[C]{\smallbreak}
\pstart
           \noindent{}\centering{}{\pb}\textcolor{gray}{\textbf{Genève\oindex{Genf@\textbf{Genf}|pw}}}\pend
           \pstart{}{\pb}\textsc{Autriche}\oindex{Österreich@\textbf{Österreich}|pw}.\pend{}\pstart{}Herrn \textsc{Gustav Schwarzkopf}\pend{}\pstart{}Wien\oindex{Wien@\textbf{Wien}|pw}\pend{}\pstart{}\textsc{I. Tiefer Graben 23}\oindex{Wien@\textbf{Wien}!I., Innere Stadt@\textbf{I., Innere Stadt}!Tiefer Graben 23@\textbf{Tiefer Graben 23}|pw}\pend{}{\bigskip}\vspace{1em}
\pstart
           \noindent{}Lieber Guſtav, es iſt ſehr ſchön geweſen, übern Jura\oindex{Jura@\textbf{Jura}|pw}, am Genferſee\oindex{Genfer See@\textbf{Genfer See}|pw} ſind
                  wir \label{K_L04507-1v}\edtext{ſpaziren gefahren}{\lemma{\textnormal{\emph{spaziren gefahren}}}\Cendnote{\textnormal{»Tour du lac\oindex{Genfer See@\textbf{Genfer See}|pwv} (mit Schwager\pwindex{Hajek, Markus 25.\,11.\,1861 Vršac – 4.\,4.\,1941 London@\textsc{Hajek, Markus} (25.\,11.\,1861 Vršac – 4.\,4.\,1941 London)|pwv}, Schwester\pwindex{Hajek, Gisela 20.\,12.\,1867 Wien – 3.\,2.\,1953 Cambridge@\textsc{Hajek, Gisela} (20.\,12.\,1867 Wien – 3.\,2.\,1953 Cambridge)|pwv} und Hugo\pwindex{Hofmannsthal, Hugo von 1.\,2.\,1874 Wien – 15.\,7.\,1929 Rodaun@\textsc{Hofmannsthal, Hugo von} (1.\,2.\,1874 Wien – 15.\,7.\,1929 Rodaun)|pw}) – Hugo\pwindex{Hofmannsthal, Hugo von 1.\,2.\,1874 Wien – 15.\,7.\,1929 Rodaun@\textsc{Hofmannsthal, Hugo von} (1.\,2.\,1874 Wien – 15.\,7.\,1929 Rodaun)|pw} in
                        Montreux\oindex{Montreux@\textbf{Montreux}|pw} Abschied«, siehe A. S.: \emph{Tagebuch}, 17. 8. 1898. Dass die Karte
                  in Ouchy\oindex{Ouchy@\textbf{Ouchy}|pwk} aufgegeben wurde, deutet darauf
                  hin, dass die Radtour dort endete und die Zielorte Montreux\oindex{Montreux@\textbf{Montreux}|pwk} und abends Genf\oindex{Genf@\textbf{Genf}|pwk} mit dem
                  Damfschiff angesteuert wurden.}}}\label{K_L04507-1}, nur wenige Gabeln ſind gebrochen, wenige
               Pneumatiks geplatzt; jetzt geht Hugo\pwindex{Hofmannsthal, Hugo von 1.\,2.\,1874 Wien – 15.\,7.\,1929 Rodaun@\textsc{Hofmannsthal, Hugo von} (1.\,2.\,1874 Wien – 15.\,7.\,1929 Rodaun)|pw} nach dem
               Süden, ich an den Vierwaldſtätterſee\oindex{Vierwaldstättersee@\textbf{Vierwaldstättersee}|pw}. Wir
               grüßen Sie von Herzen.\pend
           
\pstart
           Ihr \spacefill\mbox{Arthur}{\\[\baselineskip]}{[}hs. Hofmannsthal:{]} und \spacefill\mbox{Hugo}\pend
           \leftskip=0em{}\selectlanguage{ngerman}\vspace{1em}
\pstart
           {[}hs. G. Hajek:{]} \label{T_L04507-1v}\edtext{Beſten Gruß{\\[\baselineskip]}\spacefill\mbox{Giſela Hajek}{\\[\baselineskip]}\spacefill\mbox{{[}hs. M. Hajek:{]} Dr. M. Hajek}}{\lemma{\textnormal{\emph{Besten … Hajek}}}\Cendnote{\textnormal{Die Grüße der Hajeks\pwindex{Hajek, Gisela 20.\,12.\,1867 Wien – 3.\,2.\,1953 Cambridge@\textsc{Hajek, Gisela} (20.\,12.\,1867 Wien – 3.\,2.\,1953 Cambridge)|pwk}\pwindex{Hajek, Markus 25.\,11.\,1861 Vršac – 4.\,4.\,1941 London@\textsc{Hajek, Markus} (25.\,11.\,1861 Vršac – 4.\,4.\,1941 London)|pwk} befinden sich am rechten oberen Rand der Karte.}}}\label{T_L04507-1}\pend
           \leftskip=0em{}\selectlanguage{ngerman}\endnumbering\briefempfaengerindex{Schwarzkopf, Gustav@\textsc{Schwarzkopf, Gustav}!zzzHajek, Markus@\emph{von Markus Hajek}!1898-08-171@{17. 8. 1898}|)be}\briefempfaengerindex{Schwarzkopf, Gustav@\textsc{Schwarzkopf, Gustav}!zzzHajek, Gisela@\emph{von Gisela Hajek}!1898-08-171@{17. 8. 1898}|)be}\briefempfaengerindex{Schwarzkopf, Gustav@\textsc{Schwarzkopf, Gustav}!zzzHofmannsthal, Hugo von@\emph{von Hugo von Hofmannsthal}!1898-08-171@{17. 8. 1898}|)be}\briefempfaengerindex{Schwarzkopf, Gustav@\textsc{Schwarzkopf, Gustav}!zzzSchnitzler, Arthur@\emph{von Arthur Schnitzler}!1898-08-171@{17. 8. 1898}|)be}\mylabel{L04507h}
\begin{anhang}
\end{anhang}\newcommand{\dateiname}{L04507}\newcommand{\titel}{Arthur Schnitzler, Hugo von Hofmannsthal, Gisela und Markus Hajek an Gustav Schwarzkopf, 17. 8. 1898}\newcommand{\editorInnen}{Herausgegeben von Jahnke, SelmaMüller, Martin Anton}\input{../tex-inputs/latex-leseansicht-abspann}
